\chapter{Modeling LTI Systems}
\label{ch:linsys}
Before analyzing any specific circuit, this chapter reviews how a
\textbf{linear time-invariant (LTI)} system---one whose behavior obeys
superposition and does not change over time---is \emph{modeled}: the differential equation, the state-space model, and the transfer
function, together with the time responses they predict. These are not competing
descriptions but three ways of writing the same system down, and this chapter's job
is to build all three and move between them. Its companion, \cref{ch:splane}, then
\emph{reads} those models in the frequency and complex-plane view (Bode plots,
poles, zeros, root locus), and \cref{ch:feedback} closes the loop around them.
This chapter stops at two energy stores, which is not a limitation so much as a
choice of alphabet: \cref{ch:highorder} shows that every higher-order system is a
cascade of the first- and second-order pieces built here, so nothing beyond them
needs to be learned.
Everything here is generic---the same mathematics whether the system is a
mechanical suspension, a chemical reactor, or an electrical filter---and Part~III
then shows that each amateur-radio circuit is one of these systems in electrical
form.

\section{Four Equivalent Views of a Linear System}
An LTI system can be written four equivalent ways, and moving
between them is the central skill of this book:
\begin{enumerate}
\item a \textbf{differential equation} relating output \(y(t)\) to input \(u(t)\);
\item a \textbf{state-space model} \(\dot{\state}=A\state+Bu,\; y=C\state+Du\);
\item a \textbf{transfer function} \(G(s)=Y(s)/U(s)\) in the Laplace variable
\(s\);
\item a \textbf{frequency response} \(G(\jj\omega)\), the transfer function on the
imaginary axis.
\end{enumerate}
This chapter builds the first three and the time response; \cref{ch:splane} takes
up the fourth.

\begin{physicalbox}
\textbf{Why four?} A model is a choice of coordinates, not a claim about reality.
The differential equation is closest to the physics, the state-space model exposes
the internal variables and generalizes to many inputs and outputs, the transfer
function makes cascading and feedback into algebra, and the frequency response is
what an instrument actually measures. Fluency is being able to start from whichever
one you are handed.
\end{physicalbox}

\begin{controlsbox}
\textbf{Refresher: how the four views connect.} Taking the Laplace transform of
the state-space model (zero initial conditions) and eliminating \(\state\) gives
\[
G(s)=C\,(sI-A)^{-1}B+D .
\]
The \emph{poles} of \(G\) are the eigenvalues of \(A\); they are the system's
natural modes. Setting \(s=\jj\omega\) samples the response to a steady sinusoid
of frequency \(\omega\). So the eigenvalues you compute for stability, the poles
you read on an s-plane, and the peaks you see on a Bode plot are all the same
numbers viewed differently.
\end{controlsbox}

\section{Laplace Transforms and Block Diagrams}
The explicit bridge from a differential equation to a transfer function is the
Laplace transform. It is defined by the integral
\[
\mathcal{L}\{f\}(s)=F(s)=\int_0^\infty f(t)\,\ee^{-st}\,\dd t,
\]
which converts a function of time \(f(t)\) into a function of the complex
frequency \(s=\sigma+\jj\omega\). We rarely evaluate this integral directly;
instead, three of its properties are most useful. With zero initial
conditions, differentiation becomes multiplication by \(s\) and integration
becomes division:
\[
\mathcal{L}\{\dot f\}=sF(s),\qquad
\mathcal{L}\{\ddot f\}=s^2F(s),\qquad
\mathcal{L}\Bigl\{\textstyle\int_0^t f\,\dd t'\Bigr\}=\frac{F(s)}{s}.
\]

\begin{controlsbox}
\textbf{Refresher: Laplace generalizes phasors.} The phasor substitution
\(\dd/\dd t\to\jj\omega\) is the Laplace rule \(\dd/\dd t\to s\) evaluated on the
imaginary axis \(s=\jj\omega\). The Laplace variable \(s=\sigma+\jj\omega\) simply
adds the growth/decay coordinate \(\sigma\), which is what lets it describe
transients and stability, not just steady sinusoids.
\end{controlsbox}

Applying the derivative rule to a linear ordinary differential equation (ODE)
turns calculus into algebra: replace every \(\dd/\dd t\) by \(s\), then solve for
\(Y/U\). The two sections that follow do exactly this, once for a first-order
system and once for a second-order one. It is also the route the modeling chapters
of Part~III use, once they have written a circuit's differential equation from the
fundamentals of physics.

A transfer function is drawn as a block with an input and an output arrow, and
larger systems are built by wiring blocks together. Three rules cover almost
everything, and \cref{fig:blocks} shows all three:
\begin{description}
\item[(a) Cascade --- multiply] blocks in series multiply, \(G=G_1G_2\), because
the first block's output is the second block's input.
\item[(b) Parallel --- add] blocks driven by the same input whose outputs are
summed add, \(G=G_1+G_2\).
\item[(c) Negative feedback --- divide] a loop with forward path \(G\) and feedback
path \(H\) collapses to
\[
G_{\mathrm{cl}}=\frac{G}{1+GH}.
\]
\end{description}
\Cref{ch:feedback} takes the third rule apart in detail; it is the one that
changes a system's poles rather than just combining them.

\begin{figure}[htbp]
\centering
\begin{tikzpicture}[>=Latex,
  block/.style={draw,minimum width=1.05cm,minimum height=0.8cm},
  node distance=0.9cm]
% (a) cascade
\node (alab) at (-0.4,0) {(a)};
\node[right=0.25cm of alab] (ain) {\(U\)};
\node[block,right=0.5cm of ain] (g1) {\(G_1\)};
\node[block,right=0.7cm of g1] (g2) {\(G_2\)};
\node[right=0.6cm of g2] (aout) {\(Y\)};
\draw[->] (ain)--(g1); \draw[->] (g1)--(g2); \draw[->] (g2)--(aout);
\node[right=0.15cm of aout,black!60] {\small \(G=G_1G_2\)};
% (b) parallel
\node (clab) at (-0.4,-2.0) {(b)};
\node[right=0.25cm of clab] (cin) {\(U\)};
\coordinate (csplit) at ($(cin)+(0.55,0)$);
\node[block] (p1) at ($(cin)+(1.8,0.55)$) {\(G_1\)};
\node[block] (p2) at ($(cin)+(1.8,-0.55)$) {\(G_2\)};
\node[draw,circle,inner sep=1.2pt] (psum) at ($(cin)+(3.3,0)$) {\(\Sigma\)};
\node (cout) at ($(psum)+(0.75,0)$) {\(Y\)};
\draw (cin)--(csplit);
\draw[->] (csplit)|-(p1.west);
\draw[->] (csplit)|-(p2.west);
\draw[->] (p1.east)-|(psum.north);
\draw[->] (p2.east)-|(psum.south);
\draw[->] (psum)--(cout);
\node[right=0.15cm of cout,black!60] {\small \(G=G_1+G_2\)};
% (c) feedback
\node (blab) at (-0.4,-4.0) {(c)};
\node[right=0.25cm of blab] (bin) {\(U\)};
\node[draw,circle,inner sep=1.2pt,right=0.5cm of bin] (sum) {\(\Sigma\)};
\node[block,right=0.8cm of sum] (g) {\(G\)};
\node[right=1.5cm of g] (bout) {\(Y\)};
\node[block,below=0.75cm of g] (h) {\(H\)};
\draw[->] (bin)--(sum);
\draw[->] (sum)--(g);
\draw[->] (g)--(bout);
\coordinate (tap) at ($(g.east)!0.55!(bout)$);
\draw[->] (tap)|-(h.east);
\draw[->] (h.west)-|(sum.south);
\node at ($(sum)+(-0.02,0.3)$) {\small\(-\)};
\node[right=0.15cm of bout,black!60] {\small \(G_{\mathrm{cl}}=\dfrac{G}{1+GH}\)};
\end{tikzpicture}
\caption{Block-diagram algebra, all three rules. (a) A cascade multiplies transfer
functions. (b) Parallel paths summed at a junction add. (c) A negative-feedback
loop with forward path \(G\) and feedback \(H\) collapses to
\(G_{\mathrm{cl}}=G/(1+GH)\).}
\label{fig:blocks}
\end{figure}

\section{Why One Sinusoid at a Time Is Enough}
\label{sec:fourier}
Almost everything that follows analyzes a system by asking what it does to a
\emph{single} sinusoid. That deserves a justification, because real signals are not
sinusoids---they are speech, keyed carriers, square waves, pulses. The
justification is one line of superposition plus one theorem, and it is worth
stating once so the rest of the book can lean on it.

\textbf{Fourier's theorem} says any reasonable periodic signal of period
\(T=2\pi/\omega_0\) is a sum of sinusoids at integer multiples of \(\omega_0\):
\[
x(t)=\sum_{n=-\infty}^{\infty}c_n\,\ee^{\jj n\omega_0 t},
\]
where \(n=1\) is the \textbf{fundamental} and \(n\ge2\) are the
\textbf{harmonics}. Now feed that into a linear system. Linearity means the
response to a sum is the sum of the responses, and we already know the response to
each individual exponential---it is multiplication by \(G\) evaluated at that
frequency. So
\[
\boxed{\;y(t)=\sum_{n}G(\jj n\omega_0)\,c_n\,\ee^{\jj n\omega_0 t}\;}
\]
\emph{The transfer function sampled at the harmonics is the entire answer.} That is
the licence for the whole method: knowing \(G(\jj\omega)\) at every frequency means
knowing what the system does to \emph{any} signal, because the signal is a sum of
frequencies and the system treats each one independently. It is also why a Bode
plot is worth drawing at all.

\begin{physicalbox}
\textbf{Why a square wave has only odd harmonics.} This is usually asserted; it
follows from a symmetry. A square wave satisfies
\[
x(t+T/2)=-x(t)
\]
---shift by half a period and it inverts. Apply that to the \(n\)th term: shifting
by \(T/2\) multiplies \(\ee^{\jj n\omega_0t}\) by \(\ee^{\jj n\pi}=(-1)^n\). For the
whole sum to flip sign, every term must flip, so \((-1)^n=-1\): only \emph{odd}
\(n\) survive. Even harmonics are forbidden by the symmetry, not merely small.
Carrying out the integrals gives
\[
x(t)=\frac{4}{\pi}\left(\sin\omega_0t+\frac{\sin3\omega_0t}{3}
+\frac{\sin5\omega_0t}{5}+\cdots\right),
\]
so the amplitudes fall as \(1/n\). A sawtooth has no such symmetry, and indeed it
contains \emph{all} harmonics.

One more reading: an envelope falling as \(1/n\) is \(-\SI{20}{\decibel}\) per
decade---the same slope as a single pole (\cref{ch:splane}), and for a related
reason. A slow spectral roll-off is the frequency-domain signature of a sharp
corner in time, and a square wave's discontinuity is as sharp as corners get.
\end{physicalbox}

\section{First-Order Systems}
The canonical first-order system has one state and one time constant \(\tau\):
\[
\tau\,\dot{y}+y=K\,u
\qquad\Longleftrightarrow\qquad
\dot{x}=-\frac{1}{\tau}x+\frac{K}{\tau}u,\quad y=x .
\]
Taking the Laplace transform of the differential equation (zero initial
conditions, \(\dd/\dd t\to s\)) turns it into algebra and gives the transfer
function directly:
\[
\tau\dot y+y=Ku
\;\xrightarrow{\;\mathcal L\;}\;
\tau s\,Y+Y=KU
\;\Rightarrow\;
(\tau s+1)\,Y=KU
\;\Rightarrow\;
G(s)=\frac{Y(s)}{U(s)}=\frac{K}{\tau s+1}.
\]
Its single pole is where the denominator vanishes, \(\tau s+1=0\); it is real and
negative,
\[
s_{\text{pole}}=-\frac{1}{\tau}.
\]

\begin{physicalbox}
\textbf{First-order systems are everywhere.} Anything with a single dominant
energy or material store relaxing toward equilibrium is first order, and the same
\(\tau\dot y+y=Ku\) describes them all:
\begin{itemize}
\item a hot object cooling toward ambient (Newton's law of cooling), \(\tau=\)
thermal mass \(\times\) resistance;
\item the concentration in one well-stirred tank flushed by fresh feed,
\(\tau=\) volume\,/\,flow rate;
\item a body's velocity settling under linear drag, \(\tau=m/c\);
\item an ideal capacitor charging through a resistor, \(\tau=RC\).
\end{itemize}
In each, one time constant \(\tau\) sets how fast the system forgets its initial
condition.
\end{physicalbox}

\subsection{The Step Response}
For a unit step input the output rises exponentially to its final value,
\[
y(t)=K\bigl(1-\ee^{-t/\tau}\bigr).
\]
The time constant \(\tau\) is the natural clock of the system:
after one \(\tau\) the response has covered \(1-\ee^{-1}\approx 63.2\%\) of the
way, and after about \(5\tau\) it is settled for practical purposes
(\cref{fig:step1}).

\begin{figure}[htbp]
\centering
\includegraphics[width=0.82\textwidth]{step_first_order}
\caption{First-order step response. One time constant reaches \(63.2\%\); five
time constants are effectively settled.}
\label{fig:step1}
\end{figure}

\begin{physicalbox}
A first-order system has a single real pole at \(-1/\tau\), and no imaginary
part, so it can never overshoot or oscillate---and therefore never \emph{rings},
ringing being the decaying oscillation of an underdamped second-order system
(\cref{sec:secondorder}). Its distance from the origin, \(1/\tau\), is the decay
rate in time and (as \cref{ch:splane} shows) the corner frequency in the frequency
domain---one number, two readings.
\end{physicalbox}

\section{Worked Example: a First-Order System}
A first-order system has \(K=5\) and \(\tau=\SI{2}{\milli\second}\). Its pole is
at \(-1/\tau=-\SI{500}{\per\second}\); the step response reaches \(63.2\%\) of
\(5\) (that is, \(3.16\)) after \SI{2}{\milli\second} and is settled after about
\SI{10}{\milli\second}.

\section{Second-Order Systems}
\label{sec:secondorder}
Two independent energy stores give a second-order system, written in
\emph{standard form} as
\[
\ddot{y}+2\zeta\omega_0\,\dot{y}+\omega_0^2\,y=\omega_0^2\,u .
\]
The two parameters are the \textbf{undamped natural frequency} \(\omega_0\) and
the dimensionless \textbf{damping ratio} \(\zeta\). Radio work usually states the
same information as the \textbf{quality factor}
\[
Q=\frac{1}{2\zeta},
\]
so a lightly damped, high-\(Q\) resonator and an underdamped control system are
the same object described in two vocabularies. Physically, \(Q\) counts stored
energy against energy lost per cycle,
\[
Q=2\pi\,\frac{\text{energy stored}}{\text{energy lost per cycle}},
\]
which is why a high-\(Q\) system rings for many cycles before its oscillation dies
away. We will use \(\zeta\) when thinking
about damping and \(Q\) when thinking about sharpness, but they carry identical
content.

Applying the Laplace transform as before, with \(\dd^2/\dd t^2\to s^2\),
\[
\ddot y+2\zeta\omega_0\dot y+\omega_0^2 y=\omega_0^2u
\;\xrightarrow{\;\mathcal L\;}\;
(s^2+2\zeta\omega_0 s+\omega_0^2)\,Y=\omega_0^2\,U
\;\Rightarrow\;
G(s)=\frac{\omega_0^2}{s^2+2\zeta\omega_0 s+\omega_0^2}.
\]
A companion-form state-space model uses \(x_1=y\), \(x_2=\dot{y}\):
\[
\dot{\state}=
\mat{0 & 1\\ -\omega_0^2 & -2\zeta\omega_0}\state
+\mat{0\\ \omega_0^2}u,
\qquad
y=\mat{1 & 0}\state .
\]
The characteristic polynomial and poles are
\[
s^2+2\zeta\omega_0 s+\omega_0^2=0,
\qquad
s=-\zeta\omega_0\pm\jj\,\omega_0\sqrt{1-\zeta^2}
=-\zeta\omega_0\pm\jj\,\omega_d,
\]
where \(\omega_d=\omega_0\sqrt{1-\zeta^2}\) is the \textbf{damped natural
frequency}. \Cref{ch:splane} plots where these poles sit and what that does to
the frequency response; here we read them in the time domain.

\begin{physicalbox}
\textbf{Second-order systems are everywhere.} Two coupled energy stores that
trade energy back and forth give oscillation and, with loss, decay---the
signature of a second-order system:
\begin{itemize}
\item a mass on a spring with a damper, or a small-angle pendulum;
\item a vehicle suspension (the classic ``does it bounce or settle?'' design);
\item an underdamped feedback or process-control loop that overshoots its
setpoint and rings before settling;
\item a resonant circuit with an inductor and a capacitor.
\end{itemize}
In all of them \(\omega_0\) sets the oscillation rate and \(\zeta\) (equivalently
\(Q\)) sets how many cycles it rings before dying out.
\end{physicalbox}

The damping ratio sorts second-order systems into three behaviors, and the pole
pair above shows why:
\begin{description}
\item[\(\zeta>1\) overdamped] the square root is real, giving two real poles: no
oscillation, sluggish approach.
\item[\(\zeta=1\) critically damped] the two poles coincide at \(-\omega_0\):
the fastest approach without overshoot.
\item[\(\zeta<1\) underdamped] the poles are a complex pair: overshoot and
ringing, for roughly \(Q\) cycles.
\end{description}

\subsection{Damping You Can See}
\Cref{fig:step2} shows step responses across the regimes. A heavily damped
system (\(\zeta\ge1\)) crawls to its final value; a lightly damped one
overshoots and rings for roughly \(Q\) cycles before settling. The peak
overshoot of an underdamped step is
\[
M_p=\exp\!\left(\frac{-\pi\zeta}{\sqrt{1-\zeta^2}}\right).
\]

\begin{figure}[htbp]
\centering
\includegraphics[width=0.82\textwidth]{step_second_order}
\caption{Second-order step responses. Low \(\zeta\) (high \(Q\)) overshoots and
rings; \(\zeta=1\) is the fastest response without overshoot; high \(\zeta\) is
sluggish.}
\label{fig:step2}
\end{figure}

\begin{physicalbox}
\textbf{Mass--spring--damper analogy.} The standard second-order system is
exactly a mass \(m\) on a spring \(k\) with a damper \(c\):
\(m\ddot{x}+c\dot{x}+kx=F\), giving \(\omega_0=\sqrt{k/m}\) and
\(\zeta=c/(2\sqrt{km})\). Underdamped is a plucked guitar string that rings;
overdamped is a screen door with a strong closer that eases shut without
slamming; critically damped is the closer tuned to shut as fast as possible
without bouncing. The RLC circuit of Part~III is this same system with
\(L\leftrightarrow m\), \(1/C\leftrightarrow k\), and \(R\leftrightarrow c\).
\end{physicalbox}

\section{Worked Example: a Second-Order System}
A second-order system has \(\omega_0=2\pi(\SI{1}{\kilo\hertz})\) and
\(\zeta=0.1\), so \(Q=1/(2\zeta)=5\). The poles are at
\[
s=-\zeta\omega_0\pm\jj\omega_0\sqrt{1-\zeta^2}
\approx-628\pm\jj\,6250\ \si{\per\second}.
\]
The step overshoot is \(M_p=\exp(-\pi\cdot0.1/\sqrt{1-0.01})\approx0.73\), a
\(73\%\) overshoot---the lightly damped, ringing behavior of the top curve in
\cref{fig:step2}. \Cref{ch:splane} shows the matching \(\SI{14}{\decibel}\)
resonant peak in the frequency domain.

\begin{exambox}
The exam never asks for \(\zeta\) or a pole plot. It asks about ringing,
settling, bandwidth, and selectivity---all statements about the poles. Learn the
generic time-domain picture here, read its frequency and s-plane face in
\cref{ch:splane}, and the RC, RL, and RLC circuits of Part~III become three
readings of the same system.
\end{exambox}
